\section{Introduction}
\label{sec:introduction}

% --- The thump test as universal folklore ------------------------------------

Every summer, shoppers worldwide perform an instinctive acoustic
experiment: they tap a watermelon with their knuckles and listen.  A
deep, resonant thump signals ripeness; a dull thud suggests an
under-ripe fruit; a hollow ring warns that the flesh has turned mealy.
This `thump test'' is remarkably consistent across cultures, yet it has
never been given a quantitative physical model that connects the tap-tone
frequency to the measurable mechanical properties of the fruit.

% --- Existing empirical work -------------------------------------------------

Yamamoto et al.~\cite{Yamamoto1980} pioneered the acoustic impulse
response method, demonstrating that the natural frequencies of intact
fruits correlate with texture parameters.  De Belie et
al.~\cite{DeBelieetal2000} extended the approach to pears, and Abbott et
al.~\cite{Abbott1997} surveyed the broader landscape of non-destructive
evaluation technologies for horticultural produce.  Chen and De
Baerdemaeker~\cite{ChenDeBaerdemaeker1993} applied modal analysis to
pineapples, linking resonant behaviour to firmness.  These studies
established empirical correlations but relied on statistical regression
rather than a mechanics-based forward model.

% --- The gap: no physics-based shell model -----------------------------------

What is missing is a first-principles model that treats the fruit as what
it physically is: a fluid-filled viscoelastic shell.  Such a model would
not merely correlate frequency with firmness but would \emph{explain} the
relationship through shell eigenfrequency theory, and---crucially---would
admit an \emph{inversion} from a single measured frequency to the elastic
modulus of the rind.

% --- Connection to prior work on biological cavities -------------------------

The mathematical framework for this problem already exists.  In our
companion papers, we modelled the human abdomen as a fluid-filled oblate
spheroidal shell and derived closed-form expressions for its modal
frequencies~\cite{Junger1986, lamb1882, Leissa1973}.  The watermelon
presents a structurally analogous problem---an oblate spheroid with a
thin elastic rind enclosing a near-incompressible fluid---albeit with
different material parameters and a far more appetising motivation.

% --- Mark Kac and the inverse problem ----------------------------------------

The title of this paper echoes Kac's celebrated question, `Can one hear
the shape of a drum?''~\cite{Kac1966}.  We pose a related but distinct
inverse problem: \emph{Can one hear the ripeness of a watermelon?}  We
show that the answer is yes, at least for the dominant flexural mode
$n = 2$, and that the inversion admits a closed-form solution.

% --- Contribution and paper outline ------------------------------------------

In this paper, we:
\begin{enumerate}
  \item derive the tap-tone eigenfrequency $f_2$ for a fluid-filled
    oblate spheroidal shell parameterised by rind thickness, elastic
    modulus, geometry, and flesh density;
  \item perform Sobol sensitivity analysis to identify the dominant
    parameter (the rind elastic modulus $E_{\mathrm{rind}}$);
  \item exploit the linearity of $f_2^2$ in $E_{\mathrm{rind}}$ to
    construct a closed-form inversion from a single tap-tone measurement;
  \item introduce a dimensionless ripeness parameter
    $\Pi_{\mathrm{ripe}}$ that collapses cultivar-dependent predictions
    onto a universal master curve; and
  \item validate the model against published resonance measurements on
    watermelons and other fruits.
\end{enumerate}

The remainder of the paper is organised as follows.
Section~\ref{sec:formulation} presents the mathematical formulation.
Section~\ref{sec:results} gives the forward model predictions, Sobol
analysis, and universal curve.  Section~\ref{sec:validation} compares
predictions with published data.  Section~\ref{sec:discussion}
discusses the physical interpretation and practical implications.
Section~\ref{sec:conclusion} concludes.
